\subsection{Examples of Formula Selecting Neural Networks}

\begin{example}[Chain Rule learning in FOL]

    Let us consider a set of rules, which are formulas
    \begin{align*}
        q_0\land \ldots\land q_{n-1} \implies q_n
    \end{align*}
    where each $q_i$ is a FOL formula depending on variables $\indvariableof{i}$ and $\indvariableof{i+1}$

    We select each $q_i$ from a set $\arbsetof{i}$ enumerated by $\indvariableof{i}$.


    Effective parametrization: Sometimes we can contract the selector cores with the candidate atom cores to increase efficiency.

    \stantikzpic{
        \begin{scope}[shift={(-2,-2)}]
            \drawtdwire{-7}{3}{$\headvariable$}
            \drawhblock{-7}{0}{2}{$\bencodingof{\exformula}$}
            \drawldwire{-5}{0}{$\selvariable$}
            \drawtdwire{-8.5}{-1}{$\indvariableof{0}$}
            \drawtdwire{-5.5}{-1}{$\indvariableof{1}$}

            \drawtextnode{-1}{0}{$=$}
        \end{scope}


        \drawtdwire{5}{5}{$\headvariable$}
        \drawhblock{5}{2}{2}{$\bencodingof{\vselectionmap}$}
        \drawldwire{7}{2}{$\selvariable$}

        \draw [->-] (1,-1) to [bend right = -10] (3.5,1);
        \drawhblock{1}{-2}{2}{$\bencodingof{\folpredicateof{0}}$}

        \draw [->-] (9,-1) to [bend right = 10] (6.5,1);
        \drawhblock{9}{-2}{2}{$\bencodingof{\folpredicateof{\seldim-1}}$}

        \draw[->-] (3,-5) to [bend right = -15] (-0.5,-3);
        \draw[->-] (3,-5) to [bend right = 15] (7.5,-3);
        \drawvariabledot{3}{-5}
        \drawtdwire{3}{-5}{$\indvariableof{0}$}

        \draw[->-] (7,-5) to [bend right = -15] (2.5,-3);
        \draw[->-] (7,-5) to [bend right = 15] (10.5,-3);
        \drawvariabledot{7}{-5}
        \drawtdwire{7}{-5}{$\indvariableof{1}$}
    }


    Using the effective parameterization formulas by the tensor network:
    \stantikzpic{

        \drawtdwire{0}{3}{}
        \drawhblock{0}{0}{2}{$\bencodingof{\formulaof{0}}$}
        \draw[-<-] (2,0) to [bend right = -30] (3,-2);
        \drawtextnode{2.75}{0.25}{\colorlabelsize $\selvariableof{0}$}

        \draw[->-] (-1.5,-4) -- (-1.5,-1);
        \drawtdwire{-1.5}{-4}{}
        \drawtextnode{-2.25}{-5}{\colorlabelsize $\indvariableof{0}$}
        \drawvariabledot{-1.5}{-4}
        \draw[->-] (-1.5,-4) to [bend right=70] (21.5,-1);
%        \drawtdwire{3}{-4}{$\indvariableof{1}$}

        \draw[->-] (3,-4) to [bend right = -20] (1.5,-1);
        \drawvariabledot{3}{-4}
        \drawtdwire{3}{-4}{$\indvariableof{1}$}
        \draw[->-] (3,-4) to [bend right = 20] (4.5,-1);


        \drawtdwire{6}{3}{}
        \drawhblock{6}{0}{2}{$\bencodingof{\formulaof{1}}$}
        \drawtextnode{8.75}{0.25}{\colorlabelsize $\selvariableof{1}$}
        \draw[-<-] (8,0) to [bend right = -30] (9,-2);

        \draw[->-] (9,-4) to [bend right = -20] (7.5,-1);
        \drawvariabledot{9}{-4}
        \drawtdwire{9}{-4}{$\indvariableof{2}$}
%        \draw[->-] (9,-4) to [bend right = 20] (10.5,-1);

        \drawtextnode{10}{-1}{$\cdots$}

        \drawtextnode{10}{-3}{$\cdots$}

        %\draw[->-] (11,-4) to [bend right = -20] (9.5,-1);
        \drawvariabledot{11}{-4}
        \drawtdwire{11}{-4}{$\indvariableof{n-1}$}
        \draw[->-] (11,-4) to [bend right = 20] (12.5,-1);

        \draw[->-] (17,-4) to [bend right = -20] (15.5,-1);
        \drawvariabledot{17}{-4}
        \drawtdwire{17}{-4}{$\indvariableof{n}$}
        \draw[->-] (17,-4) to [bend right = 20] (18.5,-1);



        \drawtdwire{14}{3}{}
        \drawhblock{14}{0}{2}{$\bencodingof{\formulaof{n-1}}$}
        \drawtextnode{16.75}{0.25}{\colorlabelsize $\selvariableof{n-1}$}
        \draw[-<-] (16,0) to [bend right = -30] (17,-2);

        \drawhblock{20}{0}{2}{$\bencodingof{\formulaof{n}}$}
        \drawtextnode{22.75}{0.25}{\colorlabelsize $\selvariableof{n}$}
        \draw[-<-] (22,0) to [bend right = -30] (23,-2);

        \drawtdwire{7}{7}{}
        \drawhblock{7}{4}{8}{$\bencodingof{\land}$}

        \drawsmalltensor{14}{12}{$\tbasis$}
        \drawtdwire{14}{11}{}
        \drawhblock{14}{8}{8}{$\bencodingof{\impbincon}$}
        \drawtdwire{20}{7}{}
        \draw (20,1) -- (20,5);







    }



\end{example}

\subsubsection{Correlation}

For example (see \figref{fig:AndSupFTDecomposition}) consider the logical neuron with single activation candidate $\{\land\}$ and two variable selectors selecting $\catorder$ atomic variables $\shortcatvariables$.
The expressivity of this network is the set of all conjunctions of the atoms
\[ \{\catvariableof{\atomenumerator} \land \catvariableof{\secatomenumerator} \wcols \atomenumerator,\secatomenumerator\in[\atomorder] \} \]


% Covariance measure
Contracting with a probability distribution, we use the tensor
\[ \hypercoreat{\selvariableof{\vselectionsymbol,0},\selvariableof{\vselectionsymbol,1}} = \contractionof{\fsnn}{\selvariableof{\vselectionsymbol,0},\selvariableof{\vselectionsymbol,1}} \]
to read of covariances as
\[ \mathrm{Cov}(\catvariableof{\atomenumerator},\catvariableof{\secatomenumerator}) = \hypercoreat{\selvariableof{\vselectionsymbol,0}=\atomenumerator,\selvariableof{\vselectionsymbol,1}=\secatomenumerator}  -
\hypercoreat{\selvariableof{\vselectionsymbol,0}=\atomenumerator,\selvariableof{\vselectionsymbol,1}=\atomenumerator}  \cdot \hypercoreat{\selvariableof{\vselectionsymbol,0}=\secatomenumerator,\selvariableof{\vselectionsymbol,1}=\secatomenumerator} \, .  \]


%	\[ \skeleton = \placeholderof{1} \land \placeholderof{2} \]
%with the candidates for each placeholder being a set of $\atomorder$ atoms.

\begin{figure}[t]
    \begin{center}
        \begin{tikzpicture}[thick, scale=0.35] % , baseline = -3.5pt

    \drawatomindices{0}{-4}
    \draw (-1,-1) rectangle (5, -3);
    \node[anchor=center] (text) at (2,-2) {$\bencodingof{\atomicformulaof{\selindexof{1}} \land \atomicformulaof{\selindexof{2}}}$};

    \draw[->-] (2,-1)--(2,1) node[midway,right] {\colorlabelsize ${\headvariableof{\selindexof{1}} \land \headvariableof{\selindexof{2}}}$};

    \draw[] (7,-2.5)--(6,-2.5) node[below] {\colorlabelsize $\selvariable_{1}$};
    \drawvariabledot{6}{-2.5}
    \draw[-<-] (5,-2.5)--(6.5,-2.5); % node[midway,below] {\colorlabelsize $\selvariable_{1}$};

    \draw[] (7,-1.5)--(6,-1.5) node[above] {\colorlabelsize $\selvariable_{0}$};
    \drawvariabledot{6}{-1.5}
    \draw[-<-] (5,-1.5)--(6.5,-1.5); % node[midway,above] {\colorlabelsize $\selvariable_{0}$};

    \draw (7,-1) rectangle (9, -3);
    \node[anchor=center] (text) at (8,-2) {$\canparam$};



    \node[anchor=center] (text) at (12,-2) {${=}$};


    \begin{scope}
        [shift={(17,8)}]
        \begin{scope}
            [shift={(0,-10)}]

            \draw[->-] (5.5,5) -- (5.5,7) node[midway, right] {\colorlabelsize ${\headvariableof{\selindexof{1}} \land \headvariableof{\selindexof{2}}}$};
            \draw (1,3) rectangle (10, 5);
            \node[anchor=center] (text) at (5.5,4) {$\bencodingof{\land}$};


            % SelectorCores
            \draw[->-] (2,1) -- (2,3) node[midway, left] {\colorlabelsize ${\headvariableof{\selindexof{2}}}$};
            \draw (-1,1) rectangle (5, -1);
            \node[anchor=center] (text) at (2,0) {$\selectorcoreof{2}$};
            \draw (5,0) -- (12,0);
            \draw[-<-] (12,0) -- (13.5,0);
            \drawvariabledot{13}{0}
            \draw[] (14,0) -- (13,0) node[above] {\colorlabelsize $\selvariable_{1}$};
            \begin{scope}
                [shift={(0,-2)}]
                \draw[-<-] (0,1)--(0,-3) node[midway,left] {\colorlabelsize $\catvariableof{0}$};
                \draw[-<-] (1.5,1)--(1.5,-3) node[midway,left] {\colorlabelsize $\catvariableof{1}$};
                \node[anchor=center] (text) at (3,0) {$\cdots$};
                \draw[-<-] (4,1)--(4,-3) node[midway,right] {\colorlabelsize $\catvariableof{\atomorder\shortminus1}$};
            \end{scope}


            \draw (9,-1) -- (9,1);
            \draw[->-] (9,1) -- (9,3) node[midway, left] {\colorlabelsize ${\headvariableof{\selindexof{1}}}$};
            \draw (6,-3) rectangle (12, -1);
            \node[anchor=center] (text) at (9,-2) {$\selectorcoreof{1}$};
            \draw[-<-] (12,-2) -- (13.5,-2);
            \drawvariabledot{13}{-2}
            \draw[] (14,-2) -- (13,-2) node[above] {\colorlabelsize $\selvariable_{0}$};

            \drawatomindices{7}{-4}

            % ParameterCores
            \draw (14,1) rectangle (16, -3);
            \node[anchor=center] (text) at (15,-1) {$\canparam$};

            \begin{scope}
                [shift={(-3.5,8)}]
                \draw[fill] (7.5,-15) circle (\dotsize);
                \draw[] (7.5,-15) to[bend left=25] (3.5,-13);
                \draw[] (7.5,-15) to[bend right=25] (10.5,-13);

                \draw[fill] (9,-15.25) circle (\dotsize);
                \draw[] (9,-15.25) to[bend left=25] (5,-13);
                \draw[] (9,-15.25) to[bend right=25] (12,-13);

                \draw[fill] (11.5,-15) circle (\dotsize);
                \draw[] (11.5,-15) to[bend left=25] (7.5,-13);
                \draw[] (11.5,-15) to[bend right=25] (14.5,-13);

                \drawatomindices{7.5}{-16}

            \end{scope}

        \end{scope}
    \end{scope}

\end{tikzpicture}
    \end{center}
    \caption{Superposition of the encoded formulas ${\atomicformulaof{\selindexof{1}} \land \atomicformulaof{\selindexof{2}}}$ with weight $\canparamat{\indexedselvariableof{1},\indexedselvariableof{2}}$}
    \label{fig:AndSupFTDecomposition}
\end{figure}



\subsubsection{Conjunctive and Disjunctive Normal Forms}%\label{sec:CNFasFormulaSelection}

% Architecture
\red{
    We can represent any propositional knowledge base by the following scheme:
    Literal selecting neurons are logical neurons with connective identity/negation (selecting positive/negative literal) and selecting neurons select for each an atom.
    The single output neuron represents the disjunction, respectively the conjunction, combining the literal selecting neurons.
    The number of neurons defined by the maximal clause size plus one.
    Smaller clauses can be covered when adding False as a possible choice (The respective neuron has to choose the identity, otherwise the full clause will be trivial).
    This architecture will be discussed in more detail in \charef{cha:approximation} as $\cpformat$ selecting networks.
% Parameter
    The parameter core is in the basis $\cpformat$ format and each slice selects a clause of the knowledge base.
    In combination with polynomial decompositions, which will be provided in \charef{cha:networkRepresentation}, one can exploit this architecture to find sparse formula decompositions.
%When taking the slice values to infinity (e.g. by an annealing procedure), the represented member of the exponential family converges to the uniform distribution of the models of the knowledge base.
}

% Representation by selection tensor networks
\begin{remark}[Minterms and Maxterms]
    All minterms and maxterms can be represented by a two layer selection tensor networks without variable selection in two layers.
    The bottom layer has an $\lnot/\mathrm{Id}$ connective selection neuron to each atom and the upper layer consists of a single $\atomorder$ary conjunction.
\end{remark}




\subsection{Extension to Variables of Larger Dimension}

While we here restricted on boolean variables, formula selecting networks can be extended to variables of larger cardinality.
\begin{itemize}
    \item Connective selecting tensors: Can encode arbitrary functions $h_{\selindex}$ of discrete variables, but need $\catvariableof{\cselectionmap}$ to be an enumeration of the states, in particular to be of dimension
    \[ \catdimof{\cselectionmap} = \cardof{ \cup_{\selindexin} \imageof{h_{\selindex}} } \, . \]
    \item Variable selecting tensors can be understood as specific cases of connective selecting tensors and can thus also be generalized in a straight forward manner by
    \[ \catdimof{\cselectionmap} = \cardof{ \cup_{\selindexin} \imageof{h_{\selindex}} } \, .  \]
    \item State selecting tensors are directly defined for larger dimensions
\end{itemize}


An example of such a more generic usage is a discretization scheme for continuous neurons.

\begin{example}[Discretization of a continuous neuron]
    Let there be a neuron by a map of weight vectors and input vectors to $\rr$, that is
    \[ \lneuron( w, x)\defcols \rr^{\catorder} \times \rr^{\catorder} \rightarrow \rr \, .\]
%	When $w \in \arbsetof{weight}\subset \rr^{\catorder}$ and $x \in \arbsetof{x}\subset \rr^{\catorder}$ have
    We restrict the weights to a subset $\arbsetof{weight}\subset\rr^{\catorder}$ and the input vectors to $\arbsetof{x}\subset\rr^{\catorder}$,
    If follows that
    \[ \cardof{\imageof{\restrictionofto{\lneuron}{\arbsetof{weight}\times\arbsetof{x}}}} \leq \cardof{\arbsetof{weight}} \cdot \cardof{\arbsetof{x}} \, . \]
    To discretize the neuron, we use the subset encoding scheme of \defref{def:subsetEncoding} and define enumeration variables $\indvariableof{weight}$, $\indvariableof{x}$ and $\indvariableof{\lneuron}$ enumerating $\arbsetof{weight}$, $\arbsetof{x}$ and $\imageof{\restrictionofto{\lneuron}{\arbsetof{weight}\times\arbsetof{x}}}$, which are accompanied by respective index interpretation functions.
    Then the basis encoding of the discretized neuron is
    \begin{align*}
        \bsencodingofat{\lneuron}{\indvariableof{\lneuron},\indvariableof{weight},\indvariableof{x}} \, .
        = \sum_{\indindexofin{weight},\indindexofin{x}}
        \onehotmapof{\invindexinterpretationofat{\lneuron(\indexinterpretationofat{weight}{\indindexof{weight}},\indexinterpretationofat{x}{\indindexof{x}})}{\lneuron}}{\indvariableof{\lneuron}}
        \otimes \onehotmapofat{\indindexof{weight}}{\indvariableof{weight}}
        \otimes \onehotmapofat{\indindexof{x}}{\indvariableof{x}} \, .
    \end{align*}
    If the neuron is of the form
    \[ \lneuron(w,x) = \psi(\sum_i w_i \cdot x_i)\]
    a decomposition into multiplication at each coordinate and summation of the results, with basis encodings for each, can be done.
    \red{Here the index interpretation variables are split into a selection enumerated by $i$ and each variable gets assigned to single cores in the decomposition.}
\end{example}
